\documentclass[12pt]{article}



\newcommand{\dzbIn}[2]{\expandafter\newcommand\csname #1\endcsname{#2}} 
\newcommand{\dzbOut}[1]{\csname #1\endcsname} 

 \dzbIn{dateWeek1}{Feb 08}     
 \dzbIn{dateWeek2}{Feb 15}     
 \dzbIn{dateWeek3}{Feb 22}     
 \dzbIn{dateWeek4}{Feb 29} 
 \dzbIn{dateWeek5}{Mar 07}     
 \dzbIn{dateWeek6}{Mar 14} 
 \dzbIn{dateWeek7}{Mar 28}     
 \dzbIn{dateWeek8}{Apr 04}     
 \dzbIn{dateWeek9}{Apr 11}     
 \dzbIn{dateWeek10}{Apr 18}     
 \dzbIn{dateWeek11}{Apr 25}     
 \dzbIn{dateWeek12}{May 02}     
 \dzbIn{dateWeek13}{May 07}
 \dzbIn{midterm1Date}{Mar 14} 
 \dzbIn{midterm2Date}{Apr 23}  
 \dzbIn{finalDate}{May ??}

% Easy way to fix the margins with a package
\textwidth=7in
\textheight=9.5in
\topmargin=-1in
\headheight=0in
\headsep=.5in
\hoffset -.85in


\usepackage{amssymb, latexsym, amsmath, verbatim, multicol}
\newcommand{\N}{{\mathbb N}}
\DeclareMathOperator{\id}{id}
\newcommand{\Z}{{\mathbb Z}}
\newcommand{\C}{{\mathbb C}}
\newcommand{\R}{{\mathbb R}}
\newcommand{\Q}{{\mathbb Q}}

\setcounter{secnumdepth}{0}

\usepackage{hyperref}
\hypersetup{
 colorlinks=true,
 linkcolor=blue,
 filecolor=magenta,
 urlcolor=cyan,
}
\pagestyle{empty}

\begin{document}
\noindent\rule{18cm}{0.4pt}
\begin{center}
 {\bf MATH 220, Mathematical Reasoning and Proof
}

{\bf
   % Section 001 Wed 11:20 - 12:35 
 Section 02 TuTh 1 - 1:50
}
 \end{center}
\noindent\rule{18cm}{0.4pt}

% \renewcommand{\arraystretch}{2}

\vskip.25in
\begin{center}
  \noindent\textbf{Really rough outline of course}\\
  Last updated: \today
\end{center}
\noindent\rule{18cm}{0.4pt}
\tableofcontents

\newpage
\section[1 (\csname dateWeek1\endcsname): Introduction to course. Mathematical reasoning]{Topic  1: Introduction to course. Mathematical reasoning}

\begin{itemize}
\item Syllabus; call me DZB. Office hours, $\$\$$.
\item Introduction. How do we know what we know? Rigor and proof. Prose, ``writing requirement". 
\item Goal of the course: learn to write proofs; communicate; solve problems; be precise (``say exactly what you mean"); prepare for 300+ level courses. Create/discover new math; learn to ``generalize". Structured thinking.
\item Why rigor/proof? You know how to use procedures to do calculations; how do you know these procedures give a correct answer? 
\item Procedures $\to$ calculations $\to$ how to interpret results $\to$ why are results correct $\to$ very ``type" of input; does procedure still work? Edge cases?

\item Foundations. Inductive vs deductive reasoning. (Emphrical/data/observation vs proof/logic/symbolic manipulation.)
\item Example: even plus even is even. Examples. Authority. Define even. Examples; 6 is even; 7 is not; -8? 0? (Show magic card; do some examples of latex.) Do proof. 
\item Talk about structure of proofs. This is a "direct proof". Four steps: find hypotheses and assume them. Write out what they mean. Manipulate hypotheses. Conclude. Assume expand manipulate conclude.
\item Do proof of $2 \mid a$ and $2 \mid b$ implies that $2 \mid a + b$?
\item Common sets (Z,R,C,N,E); elements, inclusion.

\item Statements; combining statements; negations; implications; statement forms; truth tables; identities; `boolean algebra'.

\item []       \vspace{0.1cm}\hrule\vspace{-0.4cm}
\item Thursday
\item []       \vspace{0.1cm}\hrule\vspace{-0.4cm}  
\item \textbf{Negations}. Define negation. Do Handout 1, problem 3
\item \textbf{Do a few truth tables, slowly}. Do $\neg (P \vee Q)$, $(P \vee Q) \wedge R$, and $\neg \neg P$. Include $P \vee Q$ in the table.
\item  Boolean identities worksheet. Discuss and write out 1,2,3,6,7,8.
\item Do a few more identities.
\item Talk about implications again (if-then vs ``Whenever'', and how to disprove), converse, contrapositive, then do \#7 (if behind, just 1 and 5). Then do \#9 (not 3,4).

\item If extra time, do so many negations. Then do some more proofs with divisibility.
\end{itemize}


\newpage
\section[2 (\csname dateWeek2\endcsname): Divisibility problems.]{Topic  2: ``Direct'' proofs, proof by cases, and divisibility problems. }


\begin{itemize}
\item Discuss  ``depth search'' when doing proofs. Start each sentence with a word, end with punctuation. 
% \item Discuss \href{https://www.maa.org/math-competitions/putnam-competition}{Putnam}, \href{https://sites.math.northwestern.edu/~mlerma/problem_solving/putnam/easy_putnam_problems.pdf}{easy problems}, Poonen, big mouth theorem.
  % \href{http://www.math.emory.edu/~ckeyes3/drp.html}{DRP}
\item Comment on suggested problems. (Don't do them all.) Same for handouts.
\item Do a few hard nested negations (like continuity).
\item Definition of divisibility. Examples. What does it mean to prove that $a \mid b$? ``Exhibit'' ($3 \mid 6$); ``Brute force'' ($3 \mid 7$). Division algorithm.
\item Basic properties of $|$. Transitive, additive, signs; talk about corollaries; 2 out of 3 rule. Do 29999999997.
\item Now note that ``even + even = even'' is a special case of additivitiy
\item []       \vspace{0.1cm}\hrule\vspace{-0.4cm}
\item Thursday
\item []       \vspace{0.1cm}\hrule\vspace{-0.4cm}
\item $3 | 4^n-1$; $x^n-1 = ..$, $x^n - y^n = ..$
\item $4^n + 1 = $.. if $n$ odd; so, $5 | 4^n + 1$ if $n$ is odd.
\item $5 \nmid 4^n+1$ if $n$ is even. 
\item Talk about counterexamples, and disproofs.  
  
\item proof by ``cases''; $2 \mid n(n+1)$, $2 \mid ab(a-b)$. Does $3 \mid  ab(a-b)$ for all integers $a,b$?
\item Disproofs without counterexamples: prove that for all integers $n$, $n^2 + 1$ is never divisible by 4. (``Proof by cases''.) 
\item $5 | n+2 \Rightarrow 5 | n^2 - 4$; or $n^2 + 8n + 7$ (2 proofs; factor, or write $n + 2 = 5m$ and substitute.)
% \item $n$ and $n+1$ have no common divisors other than 1
\item Define prime. Is $n^3 - 1$ ever prime? What about $n^3 + 1$?
\end{itemize}

\newpage
\section[3 (\csname dateWeek3\endcsname): Proof by contradiction. Unsolvability. Irrationality.]{Topic  3: Proof by contradiction. Unsolvability of equations. Irrationality.}




\begin{itemize}
\item Recall $n^3 - 1$ and $n^3 + 1$. When is $2^n + 1$ prime? Implies that $n$ is $2^k$. Aside on Fermat numbers.  
\item $x^2 - y^2 = 1$ has no integer solutions.
\item Explain contradiction. Template (see printout?). Revisit and highlight each step (e.g., not $P$). 
\item $\neg P \Rightarrow Q$ and $\neg Q$ implies $P$. ($A \Rightarrow B$ and $\neg B$ implies $\neg A$.)
\item Definition of rational; examples; reduced form.
\item $\sqrt{2} \not \in \Q$
\item Infinitude of primes.
\item $x + y > 20 \Rightarrow x > 10$ or $y > 10$.
\item []       \vspace{0.1cm}\hrule\vspace{-0.4cm}
\item Recall contradiction. As Booleans, and Template (assume, argue, observe, conclude).
\item War: if I send troops to $X$ then they get ambushed, so I shouldn't send troops there.
\item Chess: if I move here, then in 4 moves they have checkmate. So I shouldn't move there.
\item Show that if $a$ is rational and $b$ is irrational, then $a + b$ is irrational.
\item No uninteresting number proof.
\item If $a, b, c$ are odd integers, prove that $ax^2 + bx + c = 0$ does not have a solution $x$ such that $x$ is a rational number.  
\item $\log_2{3} \not \in \Q$
\item Prove that $x^2 - y^2 = 6$ has no integer solutions.

\item \textsf{}Prove that the equation $x^2 = 4y + 3$ has no integer solutions. Emplasize cases.

\item There is no smallest positive real number.
\item OMIT. FTA (factorization); failure.
\item Prove that $\sqrt[5]{5}$ is irrational.
\item \textbf{OMIT the rest this year}
\item Worksheet. If $2^n+1$ is prime then $n$ is even. Hint for $2^n-1$ prime implies  $n$ prime.
\item More from Handout 3?
\item ``double chess'' problem. 

\end{itemize}





\newpage
\section[4 (\csname dateWeek4\endcsname): Induction.]{Topic  4: Induction.}


\begin{itemize}
\item $\sum i$,  Gauss, ``by hand''. 
\item Explain induction, and template for induction proofs. 
\item Emphasize that $P(n)$ is an open statement; it isn't $n(n+1)/2$, for instance
\item $P(1) \wedge (P(n) \Rightarrow P(n+1)) \Rightarrow \forall n, P(n)$
\item Emphasize: $P(a)$ instead of $P(n)$; $P(n-1) \Rightarrow P(n)$ is ok. Variants: start at 2, or 4, or evens.
\item $a_i = 2a_{i-1}$ 
\item Review induction. Emphasize ``build your own variants'' (start at 2, only evens, etc).
\item $2^0 + 2^1 + \cdots + 2^{n-1}$; geometric series
\item $a_1 = 0, a_n = \sqrt{3 + 2a_{n-1}} < 3$
\item []       \vspace{0.1cm}\hrule\vspace{-0.4cm}
  
\item Try to get here. 
\item []       \vspace{0.1cm}\hrule\vspace{-0.4cm}    
\item Let $n \in \mathbf{Z}_{\geq 0}$. Prove that $\displaystyle{\sum}_{i = 1}^n i^2 = \frac{n(n+1)(2n+1) }{6}$. 
\item They asked so many questions

\item Prove that, for all $n > 1$, $\frac{1}{1\cdot 2} + \frac{1}{2 \cdot 3} + \cdots + \frac{1}{(n-1)n} = \frac{n-1}{n}$
\item Fibonacci's; example, defn, lucas
\item warning: $F_{a+2} = F_{a+1} + F_{a}$, or $F_{a+1} = F_{a} + F_{a-1}$, or $F_{a} = F_{a-1} + F_{a-2}$,
\item or even $F_{2a} = F_{2a-1} + F_{2a-2}$,
\item $F_1 + F_3 + \cdots + F_{2n-1} = F_{2n}$
\item for all $n$, $F_n < 2^n$
\item $F_{n-1}F_{n+1} = F_n^2 + (-1)^{n}$
\item []       \vspace{0.1cm}\hrule\vspace{-0.4cm}  
\item Skip these this year
\item $3^a$ odd for all $a$
\item some awkward problems: $3 | 4^n-1$ via induction  
\item ``Anytime a thing works for 2 things, it works for many'' ($a,b,$ odd implies $ab$ odd, sums, etc)
\item $n! > 2^n$ for $n \geq 4$ (this is HW, skip this year

\end{itemize}


\newpage
\section[5 (\csname dateWeek5\endcsname): Basics of set theory. Basic operations. Proofs with sets.]{Topic  5: Basics of set theory. Basic operations. Proofs with sets.}


\begin{itemize}
\item Shpeal about sets (containers, some examples, anything can be a set). Definition, element of.
\item $\{1,2,3,4,5\}$, $\{1,\sqrt{2}\}$, use $\dots$ to denote a pattern. Common sets $\N, \Z$ etc2
\item Anything can be a set. $2\Z$, $d\Z$, $a\Z + b$, Fun, $\Q[x]$,  $[a,b]$ etc
\item Definition: empty set ($x \in $ empty set is always false)
\item Sets can be elements of sets $\{\{1\}, \{2\}, \{1,2\}\}$
\item Formulas to create new sets. 
\item Subset. Equality. Examples. 
\item How to do proofs with sets.  Definition is an implication. What is your hypothesis? Conclusion? 
  \begin{enumerate}
  \item Assume $x \in A$.
  \item Write out what $x \in A$ means.
  \item Argue. Or, calculate. ???
  \item Conclude that $x \in B$. (B has some definition. In step 3, verify this definition.)
  \end{enumerate}
\item To do disproofs... negate definition.
\item Do worksheet. Do at least one proof. Remember not to do the homework in class.
\item Do $6\Z \subseteq 2\Z$. Converse.
\item Do $\{4^n - 1 : n \in \Z_{\geq 0}\}$ vs $3\Z$.
\item Prove transitivity of $\subset$. 
\item Prove $\emptyset \subseteq S$ for all $S$



\item Define: union, intersection, complement, other complement, draw pictures, example with 1,2,3, and 2,3,4, 1..10
\item $6\Z$ vs $2\Z \cap 3\Z$ and $2\Z \cup 3\Z$ 
\item Propositions.
  \begin{enumerate}
  \item $A \subseteq A \cup B$
  \item $A \cup \emptyset = A$
  \item $A - B = A \cap \overline{B}$
  \item (Didn't get to the proof of this one). $A \subseteq B$ iff $A \cup B = B$.
  \end{enumerate}
\end{itemize}
% \begin{enumerate}  
% \item Handout 8 (similar problems, and propositions with intersections)
% \item $A \cap B \cap C$
% \item $(A \cap B ) \cup C = A \cap (B  \cup C)$
% \item $(A \cap B ) \cup C = (A \cup C) \cap (B  \cup C)$  
% \item $B \subseteq C \Rightarrow A-C \subseteq A-B$
% \item $B \subseteq C \Rightarrow A-C \subseteq B-C$
% \item $A \cap B \subseteq C \Rightarrow (A-C)\cap B = \emptyset$
%  \item Prove that $(A \subseteq C) \wedge (B \subseteq C) \Rightarrow A \cap B \subseteq C$.
%  \item State the converse of the previous problem. Prove or disprove it. 
% \item For all sets $A,B,C$, if $A,B$ are subsets of $C$, then $(C-A) - B = C-(A-B)$.
% \end{enumerate}


\newpage
\section[6 (\csname dateWeek6\endcsname): More sets. DeMorgan's laws. Cartesian Products. Power sets.]{Topic  6: More proofs with sets. DeMorgan's laws. Cartesian Products. Power sets}


\begin{itemize}
\item $\overline{A \cup B} = \overline{A} \cap \overline{B}$.
\item $A \cup B \subset C \Rightarrow (A - C) \cap B = \emptyset$.
% \item $A \cup B - A \cap B \subseteq (A - B) \cup (B - A)$. (Converse is HW.)

\item Sets can be elements of sets. Examples. $\{A \subseteq \mathbb{Z} \,|\, 1 \in A\}$. See notes-11.
\item Power sets. Examples. $\Z, \R$, $\{1\}$, $\{0,1\}$, $\{0,1,2\}$; $\emptyset$. 
\item $\emptyset, A \in P(A)$.
\item []       \vspace{0.1cm}\hrule\vspace{-0.4cm}    
% \item $A \subseteq B \Rightarrow P(A) \subseteq P(B)$. Converse.
\item True and false worksheet. (Do some of the HW, leave others \textbf{DZB})
 \item $P(A) \cup P(B) \subseteq P(A \cup B)$.
 \item $P(A \cup B) \subseteq P(A) \cup P(B)$.
 \item $P(A) \times P(B) \subseteq P(A \times B)$.
 \item $A = B$ if and only if $P(A) = P(B)$.
\item  Which of the following statements are true?
\item $\overline{A \cap B} = \overline{A} \cup \overline{B}$.
\item $A \cup B - A \cap B \subseteq (A - B) \cup (B - A)$. (Converse is HW.)
\item \textbf{Find more identities!}

  
  \begin{enumerate}
 \item $\{0\} \subseteq P(\mathbb{Z})$;\vspace{0.3cm}    
 \item $\{0\} \in P(\mathbb{Z})$;\vspace{0.3cm}
 \item $\mathbb{Z} \in P(\mathbb{Z})$;\vspace{0.3cm}
 \item $d\mathbb{Z} \in P(\mathbb{Z})$;\vspace{0.3cm}   
 \item $\{1,\{1\}\} \subseteq P(\mathbb{Z})$.\\ 
 \item $\{\{0,1\},\{1\}\} \subseteq P(\mathbb{Z})$;\vspace{0.3cm}
 \item $\emptyset \in P(\mathbb{Z})$;\vspace{0.3cm}
 \item $\{\emptyset\} \in P(\mathbb{Z})$. \vspace{0.3cm}
 \item $\{\emptyset\} \subseteq P(\mathbb{Z})$;\vspace{0.3cm}

 \end{enumerate}
\item []       \vspace{0.1cm}\hrule\vspace{-0.4cm}     
\item Products (emphasize ``vectors''). Examples (usual, finite, and mixed, like $\{1,-1\} \times \Z$, or $\R \times Fun$.
\item  $A \subseteq B$ and $C \subseteq D$ implies $A \times C \subseteq B \times D$.
\item $\emptyset \times A = \emptyset.$
 
\item Then start on functions.
\end{itemize}



% \newpage
% \section[  (???)]{Midterm}

% \section[  (March \csname dateWeek1\endcsname) Spring Break ]{Spring break}

% \noindent\textbf{NO CLASS March 8, 10; spring break}



\newpage
\section[7 (\csname dateWeek7\endcsname): Introduction to functions; images and surjectivity.]{Topic  7: Introduction to functions; images and surjectivity}

\textbf{Next year, add in some intervals around a point for more prep for real analysis}
\begin{itemize}
\item What is a function, pictures, ``rule''
\item $(x,y) \mapsto (x+y,xy)$. Can't graph, $f(0,0), f(1,1)$, ``unambigious''. Picture (plane to plane)
\item A and B can be anything. Let go of formulas.
\item Finite sets. ``Football play''
\item ``x is wearing glasses''; Collatz function; indicator function of $\Q$. 
\item Power set (S mapsto $S \cup {1}$, or intersect $2\Z$).
\item $\R \to P(\R)$, $x \mapsto (x,\infty)$ or $\{x\}$
\item Recursion: $n \mapsto n\cdot f(n-1)$; $(n,m) \mapsto f((n,m-1),m-1)$, $(n,0) \mapsto n$. 
\item Define Fun(A,B), $x \mapsto x^2$ is an element
\item First side of worksheet.
\end{itemize}
\begin{itemize}
\item Definition of image, surjective; $f(W) \subseteq B$; $f(W)$ is a set. To prove surj, only need $f(A) = B$.
\item Put this in a box: $a \in A \Rightarrow f(a) \in f(A)$. 
\item $b \in f(A)$ iff.. negate.
\item 1,2,3 to 4,5
\item $x^2$ is not surjecive
\item $2x+1\colon \R \to \R$ is, $\Z \to \Z$ isn't
\item cos isn't; modify codomain.
\item Second side of worksheet.
\item $X \subseteq Y \Rightarrow f(X) \subseteq f(Y)$; converse
\item $f(X\cap Y) = f(X) \cap f(Y)$
\item Monday: Common functions and images.
\item   Let $f\colon A \to B$ be a function and let $X,Y \subseteq A$. Prove
  or disprove each of the following:
  \begin{enumerate}
  \item $X \subseteq Y \Rightarrow f(X) \subseteq f(Y)$.
  \item $X \subseteq Y \Leftarrow f(X) \subseteq f(Y)$.
  \item $f(X \cup Y) \subseteq f(X) \cup f(Y)$.
  \item $f(X \cup Y) \supseteq f(X) \cup f(Y)$.
  \item $f(X \cap Y) \subseteq f(X) \cap f(Y)$.
  \item $f(X \cap Y) \supseteq f(X) \cap f(Y)$.
\end{enumerate}  
\item IVT state, and one example $e^x$
\end{itemize}




\newpage
\section[8 (\csname dateWeek8\endcsname): Inverse Image (or ``Preimage'').]{Topic 8: Inverse Image (or ``Preimage'').}

\textbf{Next year, add in some intervals around a point for more prep for real analysis}
\begin{itemize}
\item Definition. Picture. Not related to inverses.
\item $x \mapsto 6x+5$, $f^{-1}(\R_{\geq 0})$.
% \item $x \mapsto x^2$, $[1,4]$, $\R$, $\R_{> 0}$
\item Box this: $a \in f^{-1}(W)$ iff $f(a) \in W$. Negate this.
\item Note: $f^{-1}(W) \subseteq $ domain
\item Big example with 6 elements on each side. 
\item $\R^2 \to R$, $(x,y) \mapsto xy$, fibers of points
\item preimage of codomain; preimage of empty
\item cos; preimage of $[0,1]$, $\R_{\geq 0}$, $[-1,1]$, $[2,3]$, random interval, write as infinite union
\item $W \subseteq V \Rightarrow f^{-1}(W) \subseteq f^{-1}(V)$; converse fails
\item \textbf{lecture 2 starts here}
\item $\Z \to \Z$, $n \mapsto 3n+1$, preimage of evens.
\item $f(n) = n$ if even, $n-1$ if odd; preimage of 1, 0, evens, odds

\item preimage of union vs union of preimage (intersection is homework)
\item $f^{-1}(f(W))$ vs $W$. ($f(f^{-1}(W))$ is hw)
\item preimage of union 
\item $e^x$, [-1,0], [0,10]; in proof, derivative is positive, so it is increasing.
\item $\Z \to \Z$, $n \mapsto n/2$ is even $2n+4$ if odd; $W = \{4\}$, $\Z_{\geq 0}$, $\Z_{\leq 0}$, odds
\item piecewise function with overlaps
\item $S \cap f^{-1}(T) \subseteq f^{-1}(f(S) \cap T)$, converse
\item find mathoverflow post with ALL such identities
\end{itemize}




\newpage
\section[9 (\csname dateWeek9\endcsname): Injectivity.]{Topic  9: Injectivity.}


\begin{itemize}
\item Definition, slogan (``distinct inputs'' give ``distinct outputs''); horizontal line test.
\item Warning about confusing with converse. Negate. Contrapositive.
\item Examples: test functions. $x^2$. ``Not inj'' is easy (give a counterexample)
\item Emphasize: to prove that a function is injective... verify the definition. Sometimes contrapositive is better.
\item Example: $(x+1)/2$; $(x-1)/(x+1)$; $(x+y,x-y,x^2 + y^2)$.
\item $x^3$ is injective; first via algebra (and quadratic formula)
\item Define increasing/decreasing; prove that this implies injective; positive derivative thus implies injective
\item $x^3 + x$; arctan.
\item $x^2$ is not injective; restrict domain to $\R_{\geq 0}$. Prove algebraically.
\item $\cos$ is not injective; restrict domain, now injective (either ``by definition'', or via derivative).
\item $\C \to \C, z \mapsto z^3$ is not injective.
\item $S \mapsto S \cap \Z$
\item Worksheet. Max out on worksheet, and proofs. 
\item $\R \times \R^* \to \R$, $x/y$ not inj
\item $(a,b) \mapsto 2^a3^b$ is injective
\item Recall $f(X) \subseteq f(Y) \Rightarrow X \subseteq Y$ is false; counterexample is not injective; add as hypothesis
\item Is a map $\R^2 \to \R$ injective? (no if continuous, yes otherwise)
\item Next time: can probably spend 1.5 days on this topic instead and start compositions
\item Let $f\colon A \to B$ be a function. Let $X,Y \subset
  A$ and let $W,V \subseteq B$. Each of the following statements are false as stated. Which
  become true if we assume that $f$ is injective or surjective? In
  each case ($f$ is injective, or $f$ is surjective), prove your
  assertion or give a counterexample.
  \begin{enumerate}
  \item $X \subseteq Y \Leftarrow f(X) \subseteq f(Y)$.
  \item (HW)$f(X \cap Y) \subseteq f(X) \cap f(Y)$.
  \item $f(X) - f(Y) \subseteq f(X - Y)$.
  % \item $f^{-1}(X \cup Y) \subseteq f^{-1}(X) \cup f^{-1}(Y)$.
  % \item $f^{-1}(X \cup Y) \subseteq f^{-1}(X) \cup f^{-1}(Y)$.
  % \item $f^{-1}(X \cap Y) \subseteq f^{-1}(X) \cap f^{-1}(Y)$.
  % \item $f^{-1}(X \cap Y) \subseteq f^{-1}(X) \cap f^{-1}(Y)$.
  \item $X \subseteq f^{-1}(f(X))$.
  \item (HW) $W \subseteq f(f^{-1}(W))$.
  \item $V \subseteq W \Leftarrow f^{-1}(V) \subseteq f^{-1}(W)$.
  \end{enumerate}
\end{itemize}





\newpage
\section[10 (\csname dateWeek10\endcsname): Composition of functions.]{Topic  10: Composition of functions.}


\begin{itemize}
\item Definition. 
\item Examples: $x^2$ composed with $x+1$; plug in values 1,2; does not commute; commuting usually doens't make sense; formula. 
\item Examples where $A$ and 5$C$ have 2 elements, and $B$ has 3. 
\item ``Test functions'' (minimal injective, minimal surjective)
\item $1/(1+x^2)$ and $e^x$
\item $\mathbb{R} \to \mathbb{Z} \to \mathbb{Z}$, floor, then \# of distinct prime factors; plug in $\pi$, decimals, etc. 
\item $\mathbb{Z}$ to powerset to powerset, $n \mapsto \{n\}$, $S \mapsto S \cap \{1\}$.
\item Composition is a map from Fun x Fun to Fun
\item Composition is associative
\item []       \vspace{0.1cm}\hrule\vspace{-0.4cm}       
\item $id_A$ and composition

\item   Let $f\colon A \to B$ and $g \colon B \to C$ be functions. Prove or
  disprove each of the following:
  \begin{enumerate}
  \item If $f$ and $g$ are injections, then $g \circ f$ is an injection.
  \item If $f$ and $g$ are surjections, then $g \circ f$ is a surjection.
  \item If $f$ and $g$ are bijections, then $g \circ f$ is a bijection.
  \item If $g \circ f$ is an injection, then $f$ and $g$ are injections.
  \item If $g \circ f$ is a surjection, then $f$ and $g$ are surjections.
  \item If $g \circ f$ is a bijection, then $f$ and $g$ are bijections.
  \item If $g \circ f$ is an injection, then $f$ is an injection.
  \item (HW) If $g \circ f$ is an injection, then $g$ is an injection.
  \item (HW) If $g \circ f$ is a surjection, then $f$ is a surjection.
  \item (HW) If $g \circ f$ is a surjection, then $g$ is a surjection.
  \item If $g \circ f$ is a bijection, then $f$ is a bijection.
  \item If $g \circ f$ is a bijection, then $g$ is a bijection.
  \item If $g \circ f$ is an injection and $g$ is a bijection, then $f$ is an
    injection.
  \end{enumerate}



\item Let $f\colon A \to B$ be a function. Let $X,Y \subset
  A$ and let $W,V \subseteq B$. Each of the following statements are false as stated. Which
  become true if we assume that $f$ is injective or surjective? In
  each case ($f$ is injective, or $f$ is surjective), prove your
  assertion or give a counterexample.
  \begin{enumerate}
  \item $X \subseteq Y \Leftarrow f(X) \subseteq f(Y)$.
  \item (HW)$f(X \cap Y) \subseteq f(X) \cap f(Y)$.
  \item $f(X) - f(Y) \subseteq f(X - Y)$.
  % \item $f^{-1}(X \cup Y) \subseteq f^{-1}(X) \cup f^{-1}(Y)$.
  % \item $f^{-1}(X \cup Y) \subseteq f^{-1}(X) \cup f^{-1}(Y)$.
  % \item $f^{-1}(X \cap Y) \subseteq f^{-1}(X) \cap f^{-1}(Y)$.
  % \item $f^{-1}(X \cap Y) \subseteq f^{-1}(X) \cap f^{-1}(Y)$.
  \item $X \subseteq f^{-1}(f(X))$.
  \item (HW) $W \subseteq f(f^{-1}(W))$.
  \item (EX, do this one) $V \subseteq W \Leftarrow f^{-1}(V) \subseteq f^{-1}(W)$.
  \end{enumerate}
  
\end{itemize}




\newpage
\section[11 (\csname dateWeek11\endcsname): Inverse functions.]{Topic  11: Inverse functions.}

\begin{itemize}
\item From now on call this \textbf{Opposite Day.}
\item Running slightly behind, so make sure to give them tools to do HW right away.
  \begin{itemize}
  \item To check if inverse, compose and compute that you get the identity.
  \item invertible iff bijective
\end{itemize}
\item Definition. Hammer in the point that $f^{-1}$ has nothing to do with $1/f$, or with preimages.
\item Examples: $x^3$ and $x^{1/3}$; $f^{-1}(8) = 2$ because $f(2) = 8$. 
\item Explain that $f(a) = b$ iff $a = f^{-1}(b)$ (apply $f^{-1}$ to both sides and use definition). Box this. Reflect axis. 
\item Example with a set of size 3. Compute a few. Draw inverse. 
\item $\mathbb{R}\to \R$, $x^5 + 4x$. No formula for inverse. Still can compute: 0, 5, 40, 100?
\item Theorem $f$ has an inverse iff $f$ is bijective. 
  \begin{enumerate}   
  \item Use theorem to help with proofs or counterexamples (e.g., $x^2$)
  \item To find $f^{-1}$ fix $f$ and solve for $g$.
  \item If you have a guess for $f^{-1}$, plug it in and verify.
  \end{enumerate}
\item $x^5 + x^3 + x$ has an inverse.

\item $(x-1)/(x+1)$. Via IVT; then solve for inverse; then verify.
\item Example: $(x,y) \mapsto (x-y, x+y)$.
\item cos has no inverse. Tan has no inverse. 
\item Explain how to modify to be invertible. (Doesn't always work (e.g., $\sin(1/x)$))
\item Things I added last time (probably skip these): $\id \circ f = f$ (already did this), inverses are unique, $f(x) = y$ iff $x = f^{-1}(y)$.

\item Handout 13, 2-5:
  
   Let $f\colon A \to B$ and $g \colon B \to C$ be functions. Prove or
  disprove each of the following:
  \begin{enumerate}
  % \item If $f$ and $g$ are invertible, then $g \circ f$ is invertible.
  \item If $g \circ f$ is invertible, then $f$ and $g$ are invertible.
  \item If $g \circ f$ is invertible, then $f$ is invertible.
  \item If $g \circ f$ is invertible, then $g$ is invertible.
  \item If $g \circ f$ is an injection and $g$ is invertible, then $f$ is an
    injection.
  \end{enumerate}
\end{itemize}



\newpage
\section[12 (\csname dateWeek12\endcsname): Relations]{Topic  12: Relations.}

\begin{itemize}
\item Notion of related, or sameness. Similar triangles. Bijective sets (same size).
\item Definition of relation. Emphasize set vs  Examples: $\leq$. 
\item Definition of equivalence relation, and equivalence class. Emphasize that equivalence classes are sets. 
\item $\leq$, $2 \mid a - b$ (parity), triangles; 
\item Mod 2: show it is an equivalence relation. Compute equivalence classes
\item Other relations: year born. Equivalence classes. Sorting. 
\item $x \sim y$ iff $x -y \in Z$.   
\item $x \sim y$ iff $x = 1$ or $y = 1$.   
\item proof that equivalence classes partition
\item []       \vspace{0.1cm}\hrule\vspace{-0.4cm}         
\item Trivial. Total.  
\item bijections
\item worksheet (make more examples)
\item Define Graph. Graph of a function gives a relation. Not generally an equivalence relation.
\item Converse (relation comes from a function if the first projection is a bijection); actually this is the definition of a function.


\end{itemize}

\newpage
\section[13 (\csname dateWeek13\endcsname): Binary operations and countability ]{Topic  13: Binary operations and countability.}

\begin{itemize}
\item Recall that ``bijection'' is an equivalence relation. 
\item What does it mean for two sets to have the same size? Bijection between them. Dancing partners analogy. (Lead's and follows)
\item What is $|S|$ of a set? Define: $|S| = n$ if there is a bijection with $\{1,\ldots,n\}$. 
\item $\infty$ (no such bijection). 
\item Definition: we say that a set is \textbf{countable} if\ldots. Explain (bijection is a way to count, or enumerate). 
\item Hilbert Hotel analogy. ($+1, \times 2; (n,a) \mapsto p_n^a$).
\item Explain with bijections: $\mathbb{Z}_{\geq 0} \to \mathbb{N}; \mathbb{N} \to evens, \mathbb{N} \to \mathbb{Z}$. 
\item Crazier: rationals. 
\item Cantor's diagonalization argument.
\item []       \vspace{0.1cm}\hrule\vspace{-0.4cm}   
\item Recall same cardinality (bijective), $\le, \ge$, Theorem that we won't prove (if inj one way and surj other way
\item Recall uncountable, countable, give examples of countable (Z, E, Q, ExE). 
\item Power sets. There does not exist a surjection $S \to P(S)$. Corollary (lots of infinities). Does there exist a universe? (Strongly inaccessible cardinal?)
\item $\text{Fun}(S, \{0,1\}) = P(S)$. But, $\text{Fun}(\mathbb{N}, \{0,1\}) = [0,1]$ (in binary, roughly). 

\item Algebraic and transcendental numbers (most numbers are not rational).
  
\end{itemize}

\end{document}

